% Terminal apparatus of the concise study: the scope note and the References
% for the sources this companion actually uses, at the loci it cites. The
% revision timestamp and the rights colophon follow from the shared
% generation record.

\runninghead{Scope and References}
\section{Appendix: Scope and Qualifications}

{\footnotesize
\textbf{Edition, formulary and occurrence.} The Mass treated is
\latin{D.~N. Iesu Christi Regis}, \latin{Dominica ultima octobris}, I classis,
of the \work{Missale Romanum}, editio typica, Typis Polyglottis Vaticanis 1962,
printed pp.~697--698, nos.~3938--3949, with its proper Preface, p.~291,
no.~1080, in the universal 1962 calendar. Its occurrence on Sunday
25~October 2026, in place of the Twenty-second Sunday after Pentecost and
without commemoration, rests on the Missal's own calendar and general rubrics,
as \texttt{research/context.md} records. The Tract and the paschal Alleluia
printed with the formulary belong to votive use on other days and are not
parts of this Mass; the Ordinary and the act of consecration that \work{Quas
primas} n.~28 prescribes for the day are outside the formulary.

\textbf{Text and English.} The Latin was read on page images of the CMAA
facsimile of the 1962 typical edition and collated with the formulary's first
publication in \work{Acta Apostolicae Sedis} 17 (1925); the collation is in
\texttt{propers/verified.md}. The scriptural English is the Douay--Rheims in
Challoner's revision. Psalms are cited in the Missal's numbering; the
Challoner text numbers the title of Ps 71 as its first verse, so that the
Introit's verse is its 71:2, and prints the Communion's two half-verses inside
one verse numbered 28:10, and the References give its own numbering. The
English of the Preface's goods is the Lasance
\work{New Roman Missal} of 1945. Neither is an approved liturgical translation
of the 1962 Missal, and no English is composed here for a liturgical text.
Patristic witnesses are quoted in the nineteenth-century English of the
editions named below, or in their Latin.

\textbf{Reception.} The three interpretations, their sources and their four senses
are developed in full in the expansive study of this same feast, which this
companion condenses; its loci, editions and bounds are recorded in
\texttt{research/scope.md} and \texttt{research/interpretations.md}. The sweep
is bounded by the corpora named there. The Collect, Secret, Preface and
Postcommunion were composed for the feast in 1925 and have no patristic
commentary as texts, and no witness comments on the Introit's
\latin{accipere \dots\ divinitatem} itself; the Fathers' reasoning on the
second psalm and on Daniel is brought to that verb by the first interpretation.
Blessed Ildefonso Schuster's \work{The Sacramentary} V is the one commentary
read on this Mass as a Mass; he treats the texts one by one and links only the
Gradual and the Alleluia. Every other link between texts is this study's,
unless a writer is named as making it. \work{Letter}~185 is cited for Augustine's
distinction between serving God as a man and as a king, with its occasion in
the laws against the Donatists. Bellarmine is read in O'Sullivan's abridged
English; Cassiodorus and Aquinas in web transcriptions of Migne's, Parma's and
Marietti's texts; Theodoret in the PG printings. Cassiodorus and Theodoret
of Cyrus are ecclesiastical writers, and St~Cyril of Alexandria and St~Cyril
of Jerusalem support particular claims; none of them is the principal source
of an interpretation.
\work{Quas primas} is cited by number in the Holy See's English.

\textbf{Chronology.} Every biblical date on page~2 is generated from the
repository's Scripture chronology corpus, whose sources are the articles of
\work{The Catholic Encyclopedia} by Aherne, Durand, Fonck, Gigot, Meistermann,
Prat, Reid, Schets, Sloet and Van den Biesen, Haydock's Douay commentary, and
the introduction to the Psalms of the \work{New American Bible, Revised
Edition}. The corpus reads the formulary from the 1962 calendar finding aid,
which stores the votive Tract as a chant of the day, so the Tract's rows
print with the rubric that limits it; the paschal Alleluia's Apoc 19:16 has
no row. No date is held for the questioning before Pilate or for the first
year of Baltasar. Gigot's later dating of Daniel and each passage's fuller
evidence are given in the expansive study's Scriptural Date and Location.

\textbf{Rights and use.} The Latin is published on 17 U.S.C. 103(b) together
with its first publication abroad in 1925, before the 1931 line recorded in
\texttt{guidance/missals.md}, applied as of 19~October 2026. The Douay--Rheims
and the Lasance English are public domain in the United States; the patristic
editions are nineteenth-century publications, and Schuster's volume was first
published in London in 1930. Quotations from the Corpus Thomisticum
transcriptions and from \work{Quas primas} are kept brief. This is a study
aid to the 1962 Missal, not an official liturgical book, translation or
critical edition; this companion's production record is
\texttt{research/production-review.md}.
\par}

\section{References}

{\footnotesize
\begin{itemize}[leftmargin=1.2em,itemsep=.15em]
\item \work{Missale Romanum ex decreto Sacrosancti Concilii Tridentini
restitutum, Summorum Pontificum cura recognitum}, editio typica (Typis
Polyglottis Vaticanis, 1962), CMAA facsimile: \work{Rubricae generales};
Sanctorale, pp.~697--698, nos.~3938--3949; \work{Praefationes}, p.~291,
no.~1080.
\item \work{Acta Apostolicae Sedis} 17 (1925), fasc.~15, pp.~666--668: Sacred
Congregation of Rites, decree \latin{Urbis et Orbis} of 12~December 1925,
with the Mass of the feast.
\item Lasance, F.~X., \work{The New Roman Missal in Latin and English}, revised
edition (New York: Benziger Brothers, 1945), pp.~813--814 (Preface of Jesus
Christ, King) and 1257--1260 (Feast of Our Lord Jesus Christ, King).
\item \work{The Holy Bible Translated from the Latin Vulgate} (Douay--Rheims,
Challoner revision), Project Gutenberg text: Ps 2:7--11; 28:1, 10; 71:1--2,
8, 11, 20; 88:1, 39--45; 109:3; Dan 1:1--6; 7:1, 14; Mt 28:18; Lk 1:33; Jn
17:1, 20; 18:28--37; Acts 13:33; 1~Cor 15:24--28; Phil 2:6--8; Col 1:2,
12--20; 4:3, 18; Heb 2:8; Apoc 1:6, 9; 5:9--12.
\item Pius~XI, \work{Quas primas}, encyclical of 11~December 1925, English on
the Holy See's website, nn.~13, 17--19, 28, 32.
\item Augustine, \work{Tractates on the Gospel of John} 115.1--4, NPNF, first
series, vol.~7 (CCEL text).
\item Augustine, \work{Enarrationes in Psalmos} (\work{Expositions on the Book
of Psalms}), NPNF, first series, vol.~8 (CCEL text): 2.6--7; 28.10--11;
71.11, 13--14.
\item Augustine, \work{Letter} 185, \work{The Correction of the Donatists},
sections 19--20, trans. J.~R. King, NPNF, first series, vol.~4 (New Advent
text).
\item Augustine, \work{De Trinitate} I.8.16--17, I.10.20--21, trans.
A.~W. Haddan (New Advent text).
\item Basil the Great, \work{Homilia in Psalmum XXVIII}, section 8, PG 29,
cols.~304C--305A (Internet Archive page images).
\item Bede, \work{Explanatio Apocalypsis}, book I, on 1:6 and 5:5--7, PL 93
(Internet Archive OCR).
\item Bellarmine, Robert, \work{Explanatio in Psalmos}, in John O'Sullivan's
English, \work{A Commentary on the Book of Psalms} (Dublin: James Duffy, 1866),
on the second, twenty-eighth and seventy-first psalms.
\item Cassiodorus, \work{Expositio Psalmorum}, on the second and seventy-first
psalms, Migne's text (Corpus Corporum transcription).
\item Chrysostom, John, \work{Homilies on John}, homily 83, NPNF, first series,
vol.~14 (New Advent text).
\item Chrysostom, John, \work{Homilies on Colossians}, homily 3, NPNF, first
series, vol.~13 (CCEL text).
\item Cyril of Alexandria, \work{Commentary on John}, book XII, on 18:36--37,
Library of the Fathers 48 (tertullian.org transcription).
\item Cyril of Jerusalem, \work{Catechetical Lectures} 15.27--32, NPNF, second
series, vol.~7 (CCEL text).
\item Hilary of Poitiers, \work{Tractatus super Psalmos} 2.30--31, 33,
42--43, ed. A.~Zingerle, CSEL 22 (1891) (Internet Archive page images).
\item Hilary of Poitiers, \work{De Trinitate} XI.39--40, NPNF, second series,
vol.~9 (CCEL text).
\item Jerome, \work{Commentarii in Danielem}, book II, on 7:13--14, PL 25
(Wikisource transcription of Migne).
\item Schuster, Ildefonso, \work{The Sacramentary (Liber Sacramentorum)},
vol.~V (London: Burns Oates \& Washbourne, 1930), pp.~202--204 (Internet
Archive scan).
\item Theodoret of Cyrus, \work{Interpretatio in Psalmos}, on the second and
seventy-first psalms, PG 80.
\item Thomas Aquinas, \work{Super Evangelium S.~Ioannis lectura}, chapter 18,
lect.~6; \work{Postilla super Psalmos}, on the second and twenty-eighth
psalms; \work{Super Epistolam ad Colossenses lectura}, chapter one,
lect.~3, 5 (Corpus Thomisticum transcriptions).
\item \work{The Catholic Encyclopedia} (New York: Robert Appleton, 1907--1912),
for the page-2 chronology: Aherne, ``Epistle to the Colossians''; Durand,
``The New Testament''; Fonck, ``Gospel of St.~John''; Gigot, ``Book of
Daniel''; Meistermann, ``Temple of Jerusalem''; Prat, ``St.~Paul''; Reid,
``Captivities of the Israelites''; Schets, ``Third and Fourth Books of
Kings''; Sloet, ``Chronology of the Kings''; Van den Biesen, ``Apocalypse.''
\item Haydock, George Leo, \work{The Holy Bible} with Haydock's commentary
(Loreto reprint, 2014), the chronological paragraph attributed to Ussher in
its commentary on the seventieth psalm.
\item \work{New American Bible, Revised Edition}, introduction to the Psalms
(USCCB website), summarized.
\end{itemize}
\par}
